% Shared preamble for HSC ICT book
% Compile with XeLaTeX or LuaLaTeX to render Bangla text correctly.
\usepackage{iftex}
\ifPDFTeX
  \usepackage[utf8]{inputenc}
  \usepackage[T1]{fontenc}
\else
  \usepackage{fontspec}
  \defaultfontfeatures{Ligatures=TeX}
  % Bengali font stack for clean book typography. Noto Serif Bengali gives
  % the best long-form reading quality; Kalpurush/Siyam Rupali remain as
  % fallbacks for machines where Noto is not installed.
  \IfFontExistsTF{Noto Serif Bengali}{%
    \setmainfont[
      Script=Bengali,
      Scale=MatchLowercase,
      BoldFont={Noto Serif Bengali},
      ItalicFont={Noto Serif Bengali},
      BoldItalicFont={Noto Serif Bengali},
      BoldFeatures={FakeBold=1.25},
      ItalicFeatures={FakeSlant=0.12},
      BoldItalicFeatures={FakeBold=1.25,FakeSlant=0.12}
    ]{Noto Serif Bengali}
  }{%
    \IfFontExistsTF{Kalpurush}{%
      \setmainfont[
        Script=Bengali,
        FakeBold=1.4,
        FakeSlant=0.18,
        Scale=MatchLowercase
      ]{Kalpurush}
    }{%
      \setmainfont[
        Script=Bengali,
        FakeBold=1.3,
        FakeSlant=0.15,
        Scale=MatchLowercase
      ]{Siyam Rupali}
    }%
  }

  \IfFontExistsTF{Noto Sans Bengali}{%
    \setsansfont[
      Script=Bengali,
      Scale=MatchLowercase,
      BoldFont={Noto Sans Bengali},
      ItalicFont={Noto Sans Bengali},
      BoldItalicFont={Noto Sans Bengali},
      BoldFeatures={FakeBold=1.2},
      ItalicFeatures={FakeSlant=0.12},
      BoldItalicFeatures={FakeBold=1.2,FakeSlant=0.12}
    ]{Noto Sans Bengali}
  }{%
    \setsansfont[
      Script=Bengali,
      FakeBold=1.3,
      Scale=MatchLowercase
    ]{Kalpurush}
  }
\fi
\usepackage{graphicx}
% Allow using [H] float specifier and better placement control
\usepackage{float}
\usepackage{xcolor}
\usepackage{enumitem}
\usepackage{newunicodechar}
\usepackage{array}
\usepackage{booktabs}
\usepackage[most]{tcolorbox}
\usepackage[unicode]{hyperref}
% Disable colored boxes around links in the PDF (remove red boxes in TOC)
\hypersetup{hidelinks, pdfborder={0 0 0}}
\usepackage{amsmath,amssymb}

% Use a system symbol font for common bullets/glyphs missing in Bengali font.
% Try the Windows 'Segoe UI Symbol' first, fall back to a safe sans-serif.
\IfFontExistsTF{Segoe UI Symbol}{%
  \newfontfamily\symbolfont{Segoe UI Symbol}%
}{%
  \newfontfamily\symbolfont{DejaVu Sans}%
}

% Map common Unicode characters to LaTeX-safe equivalents using the symbol
% font for bullets so they render even when the main Bengali font lacks them.
\newunicodechar{•}{{\symbolfont\textbullet}}
% Common Unicode -> LaTeX mappings to avoid missing-glyph warnings
\newunicodechar{→}{\ensuremath{\rightarrow}}
\newunicodechar{±}{\ensuremath{\pm}}

% Force itemize bullets to use the symbol font (avoids missing glyph boxes)
\renewcommand{\labelitemi}{{\symbolfont\textbullet}}
\renewcommand{\labelitemii}{{\symbolfont\textendash}}
\renewcommand{\labelitemiii}{{\symbolfont\textperiodcentered}}

% Define \textenglish if not provided by language packages — keep English text inline
\providecommand{\textenglish}[1]{#1}

% Helpful graphicspath entries for topic images (add more paths if needed)
\graphicspath{{chapters/Chapter01/topics/03_virtual_reality/images/}{images/}}

% Document-level settings (fonts, margins) to be standardized here.
% Table of Contents depth: 0=chapters, 1=sections, 2=subsections.
% Keep the TOC focused on chapters and main topic sections; detailed
% subsections remain in the body to avoid an overlong contents section.
\setcounter{tocdepth}{1}
\emergencystretch=2em
\hfuzz=1pt
\setlist[itemize]{leftmargin=*, itemsep=2pt, topsep=4pt}
\setlist[enumerate]{leftmargin=*, itemsep=2pt, topsep=4pt}

% Semantic boxes for the book. Writers should use these names instead of
% manual colors so the whole book can be redesigned from one place.
\newtcolorbox{learningbox}{
  breakable,
  colback=blue!4!white,
  colframe=blue!55!black,
  coltitle=white,
  title={শিখনফল},
  fonttitle=\bfseries,
  arc=2mm,
  boxrule=0.6pt
}

\newtcolorbox{definitionbox}[1]{
  breakable,
  colback=cyan!4!white,
  colframe=cyan!45!black,
  coltitle=white,
  title={#1},
  fonttitle=\bfseries,
  arc=2mm,
  boxrule=0.6pt
}

\newtcolorbox{keypointbox}[1]{
  breakable,
  colback=gray!8!white,
  colframe=gray!60!black,
  coltitle=white,
  title={#1},
  fonttitle=\bfseries,
  arc=2mm,
  boxrule=0.6pt
}

\newtcolorbox{examplebox}[1]{
  breakable,
  colback=green!5!white,
  colframe=green!45!black,
  coltitle=white,
  title={#1},
  fonttitle=\bfseries,
  arc=2mm,
  boxrule=0.6pt
}

\newtcolorbox{boardbox}[1]{
  breakable,
  colback=yellow!10!white,
  colframe=orange!70!black,
  coltitle=black,
  title={#1},
  fonttitle=\bfseries,
  arc=2mm,
  boxrule=0.6pt
}

\newtcolorbox{mistakebox}[1]{
  breakable,
  colback=red!5!white,
  colframe=red!55!black,
  coltitle=white,
  title={#1},
  fonttitle=\bfseries,
  arc=2mm,
  boxrule=0.6pt
}

\newtcolorbox{practicebox}{
  breakable,
  colback=purple!4!white,
  colframe=purple!55!black,
  coltitle=white,
  title={অনুশীলন},
  fonttitle=\bfseries,
  arc=2mm,
  boxrule=0.6pt
}

\newtcolorbox{recapbox}{
  breakable,
  colback=gray!7!white,
  colframe=gray!55!black,
  coltitle=white,
  title={১ মিনিটে রিভিশন},
  fonttitle=\bfseries,
  arc=2mm,
  boxrule=0.6pt
}

\newtcolorbox{sourcebox}{
  breakable,
  colback=white,
  colframe=gray!45!black,
  coltitle=black,
  title={উৎস-নোট},
  fonttitle=\bfseries,
  arc=1mm,
  boxrule=0.4pt
}
